% Diagram commands used by main.tex. All figures are native TikZ so the
% Overleaf package does not depend on generated image files.

% Detailed views added during the October report review. Each diagram has an
% explicit implemented/conceptual scope in its caption or surrounding text.
\tikzset{
  efbox/.style={draw=eduflexteal!75!black,rounded corners=3pt,
    fill=teal!6,align=center,
    font={\small\setstretch{1}\hyphenpenalty=10000\exhyphenpenalty=10000},inner sep=7pt},
  efderived/.style={efbox,draw=blue!65!black,fill=blue!5},
  efdecision/.style={draw,diamond,aspect=2.4,align=center,
    fill=orange!10,font=\small,inner sep=3pt},
  efarrow/.style={-{Latex[length=2mm]},thick},
  eflabel/.style={font=\scriptsize,fill=white,inner sep=2pt}
}

\newcommand{\assessmentmodeldiagram}{
\begin{figure}[H]\centering
\begin{tikzpicture}
\node[efbox,text width=3.7cm](user) at (0,0){
\textbf{users}\\PK: user\_id};
\node[efbox,text width=3.7cm](lesson) at (7,0){
\textbf{lesson}\\PK: lesson\_id\\Optional parent FK};
\node[efderived,text width=3.7cm](attempt) at (0,-3){
\textbf{quiz\_attempts}\\FK: user, lesson\\Score, pass, timestamp};
\node[efbox,text width=3.7cm](question) at (7,-3){
\textbf{questions}\\FK: lesson\_id\\Text and points};
\node[efderived,text width=3.7cm](audit) at (0,-6){
\textbf{History boundary}\\No stored submitted answers\\No question-set version};
\node[efbox,text width=3.7cm](option) at (7,-6){
\textbf{question\_options}\\FK: question\_id\\Text and correctness};
\draw[efarrow](user)--node[eflabel,left]{1 to many}(attempt);
\draw[efarrow](lesson.west)--(3.5,0)|-node[eflabel,pos=.8,above]{1 to many}(attempt.east);
\draw[efarrow](lesson)--node[eflabel,right]{1 to many}(question);
\draw[efarrow](question)--node[eflabel,right]{1 to many}(option);
\draw[efarrow,dashed](attempt)--(audit);
\draw[efarrow](lesson.east)--++(.7,0)--++(0,1.2)-|(lesson.north);
\node[eflabel,above] at (7,1.2){optional parent relationship};
\end{tikzpicture}
\caption{Assessment definitions and persisted history; dashed note identifies missing audit detail}
\label{fig:assessment-model}
\end{figure}}

\newcommand{\rewardmodeldiagram}{
\begin{figure}[H]\centering
\begin{tikzpicture}
\node[efbox,text width=3.5cm](user) at (0,0){\textbf{users}\\PK: user\_id};
\node[efbox,text width=3.5cm](stats) at (-4.5,-2.8){
\textbf{Learner statistics}\\UQ: user\_id\\XP, level, calendar dates};
\node[efderived,text width=3.5cm](award) at (0,-2.8){
\textbf{user\_badges}\\UQ: (user, badge)\\earned\_at};
\node[efderived,text width=3.5cm](progress) at (4.5,-2.8){
\textbf{Daily quest progress}\\UQ: (user, quest, date)\\Count and completed flag};
\node[efbox,text width=3.5cm](badge) at (0,-6.5){
\textbf{badges}\\Reusable condition\\Unique condition type};
\node[efbox,text width=3.5cm](quest) at (4.5,-6.5){
\textbf{daily\_quests}\\Type, target count\\XP reward};
\draw[efarrow](user)--node[eflabel,above,sloped]{1 to 0 or 1}(stats);
\draw[efarrow](user)--node[eflabel,left]{1 to many}(award);
\draw[efarrow](user)--node[eflabel,above,sloped]{1 to many}(progress);
\draw[efarrow](badge)--node[eflabel,left]{1 to many}(award);
\draw[efarrow](quest)--node[eflabel,right]{1 to many}(progress);
\end{tikzpicture}
\caption{Reward definitions, aggregate totals, and per-learner earned relationships}
\label{fig:reward-model}
\end{figure}}

\newcommand{\quizworkflowdiagram}{
\begin{figure}[H]\centering
\begin{tikzpicture}
\node[efbox,text width=12cm](owner) at (0,0){
Controller: require the authenticated owner of the submitted user ID};
\node[efbox,text width=5.7cm](mc) at (-3.4,-2.3){
\textbf{Multiple choice}\\Transaction + learner lock\\Enrolment and lesson type};
\node[efderived,text width=5.7cm](fb) at (3.4,-2.3){
\textbf{Fill-in-the-blank}\\Separate feedback use case\\No course entitlement check};
\node[efbox,text width=5.7cm](grade) at (-3.4,-5.1){
Check option and lesson scope\\Score submitted items\\Save attempt; update quests};
\node[efderived,text width=5.7cm](text) at (3.4,-5.1){
Load answer key by question ID\\Trim and ignore case\\Compare submitted text};
\node[efbox,text width=5.7cm](reward) at (-3.4,-7.9){
If first pass: 30 XP\\Complete quiz and parent\\Evaluate course benefits};
\node[efderived,text width=5.7cm](feedback) at (3.4,-7.9){
Correctness and expected text\\No attempt history\\No progress or XP};
\node[efbox,text width=5.7cm](result) at (-3.4,-10.7){
Commit or roll back\\Return score and eligible state};
\draw[efarrow](owner)--(mc);\draw[efarrow](owner)--(fb);
\draw[efarrow](mc)--(grade);\draw[efarrow](grade)--(reward);
\draw[efarrow](reward)--(result);
\draw[efarrow](fb)--(text);\draw[efarrow](text)--(feedback);
\end{tikzpicture}
\caption{Existing quiz workflows differ in scope, persistence, and rewards; grading gaps remain}
\label{fig:quiz-workflow}
\end{figure}}

\newcommand{\restoreplandiagram}{
\begin{figure}[H]\centering
\begin{tikzpicture}[node distance=.55cm]
\node[efbox,text width=11.5cm](isolate){
Isolate a restore environment; preserve the damaged state for diagnosis};
\node[efbox,text width=11.5cm,below=of isolate](db){
Restore independent PostgreSQL backup with required extension and roles};
\node[efbox,text width=11.5cm,below=of db](check){
Validate Flyway version, row relationships, progress, and reward totals};
\node[efderived,text width=5cm,below=.8cm of check,xshift=-3.2cm](cache){
Discard disposable DTO caches\\Apply explicit session invalidation policy};
\node[efderived,text width=5cm,below=.8cm of check,xshift=3.2cm](vector){
Reconcile or regenerate vectors\\Replay tool is future work};
\node[efbox,text width=11.5cm,below=1cm of cache,xshift=3.2cm](verify){
Verify media references, client journey, metrics, and alert resolution\\Record observed recovery time and recovered data boundary};
\draw[efarrow](isolate)--(db);\draw[efarrow](db)--(check);
\draw[efarrow](check.south)--(cache.north);
\draw[efarrow](check.south)--(vector.north);
\draw[efarrow](cache.south)--(verify.north);
\draw[efarrow](vector.south)--(verify.north);
\end{tikzpicture}
\caption{Proposed isolated restoration drill; backup automation and replay are not implemented}
\label{fig:restore-plan}
\end{figure}}

\newcommand{\responsibilitydiagram}{
\begin{figure}[H]\centering
\begin{tikzpicture}[node distance=.7cm]
\node[efbox,text width=12cm](client){\textbf{Interaction boundary}\\
Android navigation, loading/empty/error states, bearer credentials};
\node[efbox,text width=12cm,below=of client](rules){\textbf{Authoritative application boundary}\\
Caller ownership, enrolment, quiz evaluation, transaction coordination};
\node[efbox,text width=5.1cm,below=of rules,xshift=-3.3cm](db){
\textbf{PostgreSQL source of truth}\\Progress, XP, attempts, enrolment, constraints};
\node[efderived,text width=5.1cm,below=of rules,xshift=3.3cm](derived){
\textbf{Derived representations}\\Redis DTO caches; RabbitMQ-driven vectors};
\node[efderived,text width=12cm,below=1cm of db,xshift=3.3cm](obs){
\textbf{Observation boundary}\\Micrometer metrics, exporters, probes, dashboards, alerts};
\draw[efarrow](client)--(rules);
\draw[efarrow](rules.south)--(db.north);
\draw[efarrow](rules.south)--(derived.north);
\draw[efarrow,dashed](db)--(obs);
\draw[efarrow,dashed](derived)--(obs);
\end{tikzpicture}
\caption{Responsibility boundaries: authoritative, derived, and observed state}
\label{fig:responsibility}
\end{figure}}

\newcommand{\learningerdetaildiagram}{
\begin{figure}[H]\centering
\begin{tikzpicture}
\node[efbox,text width=3.5cm](user) at (0,0){
\textbf{users}\\PK: user\_id\\Identity and role};
\node[efbox,text width=3.5cm](course) at (7,0){
\textbf{courses}\\PK: course\_id\\Metadata, price, embedding};
\node[efderived,text width=3.5cm](enrol) at (3.5,-3.2){
\textbf{enrollments}\\FK: user\_id, course\_id\\UQ: (user, course)};
\node[efbox,text width=3.5cm](progress) at (0,-6.2){
\textbf{lesson\_progress}\\FK: user\_id, lesson\_id\\UQ: (user, lesson)};
\node[efbox,text width=3.5cm](lesson) at (7,-6.2){
\textbf{lesson}\\FK: course\_id\\Content, hierarchy, embedding};
\node[efderived,text width=3.5cm](attempt) at (3.5,-9.3){
\textbf{quiz\_attempts}\\FK: user\_id, lesson\_id\\Repeated assessment history};
\draw[efarrow](user)--node[eflabel,above,sloped]{1 to many}(enrol);
\draw[efarrow](course)--node[eflabel,above,sloped]{1 to many}(enrol);
\draw[efarrow](user)--node[eflabel,left]{1 to many}(progress);
\draw[efarrow](course)--node[eflabel,right]{1 to many}(lesson);
\draw[efarrow](lesson)--node[eflabel,above]{1 to many}(progress);
\draw[efarrow](lesson)--node[eflabel,above,sloped]{1 to many}(attempt);
\draw[efarrow](user.west)--++(-.6,0)|-node[eflabel,pos=.85,above]{1 to many}(attempt.west);
\end{tikzpicture}
\caption{Learning data model with cardinality and natural-key constraints}
\label{fig:learning-er-detail}
\end{figure}}

\newcommand{\tokenrefreshdiagram}{
\begin{figure}[H]\centering
\begin{tikzpicture}[font=\small]
\foreach \x/\name in {0/Request A,4.5/Coordinator,9/Request B}{
\node[efbox,text width=2.8cm] at (\x,0){\name};
\draw[dashed](\x,-.45)--(\x,-8.2);}
\draw[efarrow](0,-1)--node[eflabel,above]{401 with old token}(4.5,-1);
\draw[efarrow](9,-1.8)--node[eflabel,above]{401 with same old token}(4.5,-1.8);
\node[efderived,text width=5.4cm] at (4.5,-3.1){
Synchronized decision\\A refreshes through public client\\Store replacement token};
\draw[efarrow](4.5,-4.6)--node[eflabel,above]{A retries with replacement}(0,-4.6);
\node[efderived,text width=5.4cm] at (4.5,-6.2){
B enters after A\\Current token differs from sent token\\Reuse stored replacement};
\draw[efarrow](4.5,-7.8)--node[eflabel,above]{B retries; no second refresh}(9,-7.8);
\end{tikzpicture}
\caption{Implemented in-process coordination for concurrent expired-token requests}
\label{fig:token-refresh}
\end{figure}}

\newcommand{\screenstatediagram}{
\begin{figure}[H]\centering
\begin{tikzpicture}
\node[efbox,text width=2.8cm](loading) at (0,0){Loading\\Request in flight};
\node[efbox,text width=2.8cm](content) at (-4,-3.2){Content\\Render returned state};
\node[efderived,text width=2.8cm](empty) at (0,-3.2){Empty\\Offer next action};
\node[efderived,text width=2.8cm](failure) at (4,-3.2){Failure\\Explain and allow retry};
\draw[efarrow](loading.south)--(0,-1.6)-|(content.north);
\draw[efarrow](loading.south)--node[eflabel,left,pos=.8]{no data}(empty.north);
\draw[efarrow](loading.south)--(0,-1.6)-|(failure.north);
\node[eflabel,above] at (-2.3,-1.6){success with data};
\node[eflabel,above] at (2.3,-1.6){request failed};
\draw[efarrow](failure.east)--(6,-3.2)--(6,1.2)--(0,1.2)--(loading.north);
\node[eflabel,above] at (2.8,1.2){retry; preserve selected resource};
\end{tikzpicture}
\caption{Screen-state acceptance model for loaded data; device coverage remains partial}
\label{fig:screen-state}
\end{figure}}

\newcommand{\passwordresetdiagram}{
\begin{figure}[H]\centering
\resizebox{\textwidth}{!}{
\begin{tikzpicture}[font=\small]
\foreach \x/\name in {0/Learner,4/API,8/PostgreSQL,12/SMTP}{
\node[efbox,text width=2.2cm] at (\x,0){\name};
\draw[dashed](\x,-.45)--(\x,-9.5);}
\draw[efarrow](0,-1)--node[eflabel,above]{request(email)}(4,-1);
\draw[efarrow](4,-1.9)--node[eflabel,above]{store SHA-256 digest + expiry}(8,-1.9);
\draw[efarrow](4,-2.8)--node[eflabel,above]{deliver plaintext random code}(12,-2.8);
\draw[efarrow](4,-3.7)--node[eflabel,above]{generic request response}(0,-3.7);
\draw[efarrow](0,-4.6)--node[eflabel,above]{confirm(code, new password)}(4,-4.6);
\draw[efarrow](4,-5.5)--node[eflabel,above]{lock valid unused digest}(8,-5.5);
\draw[efarrow](4,-6.4)--node[eflabel,above]{consume + write BCrypt hash}(8,-6.4);
\node[efderived,text width=4.5cm] at (4,-7.3){Revoke refresh tokens in Redis\\Separate system boundary};
\draw[efarrow](4,-9.2)--node[eflabel,above]{success or invalid/expired}(0,-9.2);
\end{tikzpicture}}
\caption{Implemented recovery lifecycle when SMTP is configured; not one distributed transaction}
\label{fig:password-reset}
\end{figure}}

\newcommand{\paymentsimulationdiagram}{
\begin{figure}[H]\centering
\begin{tikzpicture}[node distance=.55cm]
\node[efbox,text width=10.5cm](owner){Validate JWT owner and simulation flag};
\node[efbox,text width=10.5cm,below=of owner](price){Load server-side course price};
\node[efbox,text width=10.5cm,below=of price](write){
Transaction: insert SUCCESS record\\ON CONFLICT DO NOTHING; accept existing successful record};
\node[efbox,text width=10.5cm,below=of write](enrol){
Ensure unique learner--course enrolment};
\node[efderived,text width=10.5cm,below=of enrol](result){
Commit and return simulation=true\\No external payment or settlement};
\draw[efarrow](owner)--(price);
\draw[efarrow](price)--(write);
\draw[efarrow](write)--(enrol);
\draw[efarrow](enrol)--(result);
\end{tikzpicture}
\caption{Retry-safe development checkout; repeated requests converge on one successful record}
\label{fig:payment-simulation}
\end{figure}}

\newcommand{\concurrenttransactiondiagram}{
\begin{figure}[H]\centering
\begin{tikzpicture}[font=\small]
\foreach \x/\name in {0/Transaction A,5/PostgreSQL,10/Transaction B}{
\node[efbox,text width=2.8cm] at (\x,0){\name};
\draw[dashed](\x,-.45)--(\x,-8);}
\draw[efarrow](0,-1)--node[eflabel,above]{lock learner row}(5,-1);
\draw[efarrow](10,-1.8)--node[eflabel,above]{request same row lock}(5,-1.8);
\node[efderived,text width=2.7cm] at (10,-3){Wait for A};
\draw[efarrow](0,-3.2)--node[eflabel,above]{new completion; add 20 XP}(5,-3.2);
\draw[efarrow](0,-4.3)--node[eflabel,above]{commit; release lock}(5,-4.3);
\draw[efarrow](5,-5.3)--node[eflabel,above]{B acquires lock}(10,-5.3);
\draw[efarrow](10,-6.3)--node[eflabel,above]{already complete; no new XP}(5,-6.3);
\draw[efarrow](10,-7.3)--node[eflabel,above]{commit existing result}(5,-7.3);
\end{tikzpicture}
\caption{Successful duplicate completion ordering for one learner under a row lock}
\label{fig:concurrent-transaction}
\end{figure}}

\newcommand{\incidentdecisiondiagram}{
\begin{figure}[H]\centering
\begin{tikzpicture}
\node[efbox,text width=11cm](alert) at (0,0){Availability alert: inspect main-listener probe and management scrape};
\node[efderived,text width=4.4cm](listener) at (-3,-2.1){
Probe fails; scrape succeeds\\Inspect listener, readiness dependencies, network};
\node[efderived,text width=4.4cm](process) at (3,-2.1){
Probe and scrape fail\\Inspect API process, JVM startup, container health};
\node[efbox,text width=11cm](recover) at (0,-4.3){
Identify cause; restore affected service or dependency\\Avoid destructive volume removal};
\node[efbox,text width=11cm](verify) at (0,-6.1){
Verify client request, readiness, scrape targets, and alert resolution};
\draw[efarrow](alert)--(listener);
\draw[efarrow](alert)--(process);
\draw[efarrow](listener)--(recover);
\draw[efarrow](process)--(recover);
\draw[efarrow](recover)--(verify);
\end{tikzpicture}
\caption{Runbook-based investigation and recovery decision model}
\label{fig:incident-decision}
\end{figure}}

\newcommand{\systemcontextdiagram}{
\begin{figure}[H]\centering
\resizebox{0.94\textwidth}{!}{
\begin{tikzpicture}[
person/.style={draw,rounded corners,fill=orange!12,align=center,minimum width=2.6cm,minimum height=1cm},
system/.style={draw,rounded corners,very thick,fill=teal!9,align=center,minimum width=4.5cm,minimum height=2cm},
external/.style={draw,dashed,rounded corners,fill=blue!5,align=center,minimum width=2.8cm,minimum height=1cm},
arr/.style={-{Latex[length=2mm]},thick}]
\node[system](platform){\textbf{EduFlex platform}\\Android client + Spring Boot API\\learning, assessment, progress, rewards};
\node[person,left=2.1cm of platform,yshift=1.1cm](learner){Learner};
\node[person,left=2.1cm of platform,yshift=-1.1cm](admin){Administrator};
\node[person,below=1.4cm of platform](operator){Operator};
\node[external,right=2cm of platform,yshift=1.5cm](supabase){Supabase Storage\\optional media};
\node[external,right=2cm of platform](gemini){Gemini API\\optional generation};
\node[external,right=2cm of platform,yshift=-1.5cm](firebase){Firebase Messaging\\optional push};
\draw[arr](learner)--(platform);
\draw[arr](admin)--(platform);
\draw[arr](operator)--(platform);
\draw[arr](platform)--(supabase);
\draw[arr](platform)--(gemini);
\draw[arr](firebase)--(platform);
\end{tikzpicture}}
\caption{System context and external actors}\label{fig:system-context}
\end{figure}}

\newcommand{\backendlayersdiagram}{
\begin{figure}[H]\centering
\resizebox{0.92\textwidth}{!}{
\begin{tikzpicture}[
layer/.style={draw,rounded corners,align=center,minimum width=11.5cm,minimum height=1cm},
arr/.style={-{Latex[length=2mm]},thick}]
\node[layer,fill=orange!10](http){\textbf{HTTP boundary}: controllers, validation, DTOs, exception mapping, OpenAPI};
\node[layer,fill=teal!9,below=.5cm of http](services){\textbf{Application rules}: authentication, courses, progress, quizzes, gamification, AI, media};
\node[layer,fill=blue!7,below=.5cm of services](repositories){\textbf{Persistence and integration}: JPA/jOOQ repositories, Redis cache, AMQP publisher and listener};
\node[layer,fill=gray!10,below=.5cm of repositories](runtime){\textbf{Runtime}: PostgreSQL/pgvector, Redis, RabbitMQ, Supabase, Gemini};
\draw[arr](http)--node[right]{validated calls}(services);
\draw[arr](services)--node[right]{queries, commands, events}(repositories);
\draw[arr](repositories)--node[right]{protocol clients}(runtime);
\end{tikzpicture}}
\caption{Backend responsibility layers}\label{fig:backend-layers}
\end{figure}}

\newcommand{\androidnavigationdiagram}{
\begin{figure}[H]\centering
\resizebox{\textwidth}{!}{
\begin{tikzpicture}[
screen/.style={draw,rounded corners,fill=teal!7,align=center,minimum width=2.2cm,minimum height=.75cm},
admin/.style={draw,rounded corners,fill=orange!10,align=center,minimum width=2.2cm,minimum height=.75cm},
arr/.style={-{Latex[length=1.7mm]},thick}]
\node[screen](auth){Login / Register};
\node[screen,right=1cm of auth](home){Home};
\node[screen,right=1cm of home](discover){Discover / Search};
\node[screen,right=1cm of discover](detail){Course Detail};
\node[screen,right=1cm of detail](lesson){Lesson Study};
\node[screen,below=.8cm of lesson](quiz){Quiz / Result};
\node[screen,below=.8cm of detail](ai){AI Summary};
\node[screen,below=.8cm of discover](cart){Cart / Payment};
\node[screen,below=.8cm of home](profile){Profile / Rewards};
\node[screen,below=.8cm of auth](learn){My Learning};
\node[admin,below=2.7cm of detail](adminhome){Admin Users};
\node[admin,left=.8cm of adminhome](admincourses){Admin Courses};
\node[admin,right=.8cm of adminhome](adminlessons){Admin Lessons};
\draw[arr](auth)--(home);\draw[arr](home)--(discover);\draw[arr](discover)--(detail);
\draw[arr](detail)--(lesson);\draw[arr](lesson)--(quiz);\draw[arr](detail)--(ai);
\draw[arr](discover)--(cart);\draw[arr](home)--(profile);\draw[arr](home)--(learn);
\draw[arr,dashed](auth.south)to[bend right=25](adminhome.west);
\draw[arr](adminhome)--(admincourses);\draw[arr](adminhome)--(adminlessons);
\end{tikzpicture}}
\caption{Android navigation and principal user journeys}\label{fig:android-navigation}
\end{figure}}

\newcommand{\deploymentdiagram}{
\begin{figure}[H]\centering
\resizebox{\textwidth}{!}{\begin{tikzpicture}
\node[efbox,text width=2.8cm](android) at (-4,2.2){Android device\\API base URL};
\node[efbox,text width=3.2cm](api) at (0,0){Spring Boot API\\HTTP 8080\\Management 8081};
\node[efbox,text width=2.8cm](pg) at (-4,-3){PostgreSQL\\pgvector extension\\Named data volume};
\node[efbox,text width=2.8cm](redis) at (0,-3){Redis\\AOF data volume};
\node[efbox,text width=2.8cm](rabbit) at (4,-3){RabbitMQ\\Management 15672\\No mounted data volume};
\node[efderived,text width=2.8cm](exporters) at (-4,-6){PostgreSQL and Redis\\exporters};
\node[efderived,text width=2.8cm](prom) at (0,-6){Prometheus 9090\\Scrape and evaluate};
\node[efderived,text width=2.8cm](blackbox) at (4,-6){Blackbox exporter\\Probe API listener};
\node[efderived,text width=2.8cm](grafana) at (-4,-8.5){Grafana 3000\\Query and visualize};
\node[efderived,text width=2.8cm](alert) at (4,-8.5){Alertmanager 9093\\Local alert receiver};
\draw[efarrow](android)--(api);
\draw[efarrow](api)--(pg);\draw[efarrow](api)--(redis);\draw[efarrow](api)--(rabbit);
\draw[efarrow](prom.west)--(exporters.east);
\draw[efarrow](exporters)--(pg);
\draw[efarrow](exporters.north east)--(redis.south west);
\draw[efarrow](prom.east)--(blackbox.west);
\draw[efarrow](prom.north east)--node[eflabel,above,sloped]{broker metrics}(rabbit.south west);
\draw[efarrow](grafana.north east)--(prom.south west);
\draw[efarrow](prom.south east)--(alert.north west);
\draw[efarrow](prom.north west)--(-6,-5)--(-6,0)--(api.west);
\node[eflabel,rotate=90] at (-6,-2.1){management scrape};
\draw[efarrow](blackbox.north east)--(6,-5)--(6,0)--(api.east);
\node[eflabel,rotate=90] at (6,-2.1){main-listener readiness probe};
\end{tikzpicture}}
\caption{Compose service dependencies: arrows follow callers; broker data storage is not mounted}
\label{fig:deployment}
\end{figure}}

\newcommand{\erdmodeldiagram}{
\begin{figure}[H]\centering
\begin{tikzpicture}
\node[efbox,text width=3.7cm](users) at (-3.4,0){\textbf{Identity}\\users};
\node[efbox,text width=3.7cm](courses) at (3.4,0){\textbf{Catalogue}\\courses};
\node[efderived,text width=4.9cm](shared) at (0,-2.5){
\textbf{Learner--course relationships}\\enrollments, course\_reviews, transactions};
\node[efbox,text width=4.2cm](state) at (-3.4,-5.2){
\textbf{Learner state}\\lesson\_progress, statistics, badge awards, daily quest progress};
\node[efbox,text width=4.2cm](content) at (3.4,-5.2){
\textbf{Course content}\\lesson hierarchy, questions, answer options, embeddings};
\node[efderived,text width=4.9cm](attempts) at (0,-8){
\textbf{Assessment history}\\quiz\_attempts connects users and lessons};
\draw[efarrow](users)--(shared);\draw[efarrow](courses)--(shared);
\draw[efarrow](users.south west)--(state.north west);
\draw[efarrow](courses.south east)--(content.north east);
\draw[efarrow](state)--(attempts);\draw[efarrow](content)--(attempts);
\draw[efarrow,dashed](content.west)--node[eflabel,above]{lesson links}(state.east);
\end{tikzpicture}
\caption{Domain relationship overview; detailed figures specify keys and cardinality}
\label{fig:erd}
\end{figure}}

\newcommand{\apiboundarydiagram}{
\begin{figure}[H]\centering
\begin{tikzpicture}
\node[efbox,text width=10cm](policy) at (0,0){Request enters security filter chain\\Evaluate route policy};
\node[efderived,text width=3.7cm](public) at (-3,-2){
Public route\\Login, register, recovery, refresh, logout, health, OpenAPI};
\node[efbox,text width=3.7cm](jwt) at (3,-2){
Protected route\\Verify access-token signature, issuer, expiry, type};
\node[efbox,text width=3.7cm](learner) at (0,-5.4){
Learner commands\\Verify caller ownership\\Enrolment where required};
\node[efderived,text width=3.7cm](admin) at (5.3,-5.4){
\texttt{ADMIN} commands\\Privileged content and user changes\\Operation-specific rules};
\draw[efarrow](policy)--(public);
\draw[efarrow](policy)--(jwt);
\draw[efarrow](jwt)--(learner);
\draw[efarrow](jwt)--(admin);
\end{tikzpicture}
\caption{API groups and authorization boundaries}\label{fig:api-boundaries}
\end{figure}}

\newcommand{\authsequencediagram}{
\begin{figure}[H]\centering
\resizebox{0.96\textwidth}{!}{
\begin{tikzpicture}[x=2.7cm,y=1.05cm,arr/.style={-{Latex[length=1.7mm]},thick}]
\foreach \x/\name in {0/Android,1/API,2/PostgreSQL,3/Redis}
  {\node[draw,rounded corners,fill=teal!7,minimum width=2.1cm] at (\x,0) {\name};
   \draw[dashed] (\x,-.35)--(\x,-7.5);}
\draw[arr](0,-1)--node[above]{1. email + password}(1,-1);
\draw[arr](1,-2)--node[above]{2. load user + BCrypt hash}(2,-2);
\draw[arr](2,-2.7)--node[above]{3. user record}(1,-2.7);
\draw[arr](1,-3.7)--node[above]{4. store refresh token ID, TTL}(3,-3.7);
\draw[arr](1,-4.5)--node[above]{5. access + refresh tokens}(0,-4.5);
\draw[arr](0,-5.5)--node[above]{6. refresh token}(1,-5.5);
\draw[arr](1,-6.2)--node[above]{7. verify not revoked}(3,-6.2);
\draw[arr](1,-7)--node[above]{8. new access token}(0,-7);
\end{tikzpicture}}
\caption{Login and access-token refresh sequence}\label{fig:auth-sequence}
\end{figure}}

\newcommand{\learnerjourneydiagram}{
\begin{figure}[H]\centering
\begin{tikzpicture}[
state/.style={draw,rounded corners,fill=teal!7,align=center,minimum width=2.3cm,minimum height=.9cm},
decision/.style={draw,diamond,aspect=2,fill=orange!9,align=center,inner sep=1.5pt},
arr/.style={-{Latex[length=1.5mm]},thick}]
\node[state](discover){Discover course};
\node[state,below=.45cm of discover](detail){Read details};
\node[decision,below=.45cm of detail](paid){Paid?};
\node[state,right=2.2cm of paid](payment){Simulate payment\\record};
\node[state,below=.55cm of paid](enrol){Create enrolment};
\node[state,below=.45cm of enrol](study){Study lesson};
\node[decision,below=.45cm of study](quiz){Quiz lesson?};
\node[state,right=2.2cm of quiz](submit){Submit answer};
\node[state,below=.55cm of quiz](complete){Save progress};
\node[state,below=.45cm of complete](reward){Award XP, quest progress,\\badge, and completion};
\draw[arr](discover)--(detail);\draw[arr](detail)--(paid);
\draw[arr](paid)--node[above]{yes}(payment);\draw[arr](payment.south)to[bend left=18](enrol.east);
\draw[arr](paid)--node[left]{no}(enrol);\draw[arr](enrol)--(study);
\draw[arr](study)--(quiz);\draw[arr](quiz)--node[above]{yes}(submit);
\draw[arr](submit.south)to[bend left=18](complete.east);\draw[arr](quiz)--node[left]{no}(complete);
\draw[arr](complete)--(reward);
\end{tikzpicture}
\caption{End-to-end learner course journey}\label{fig:learner-journey}
\end{figure}}

\newcommand{\progressstatediagram}{
\begin{figure}[H]\centering
\resizebox{0.9\textwidth}{!}{
\begin{tikzpicture}[
state/.style={draw,rounded corners,fill=blue!6,align=center,minimum width=2.6cm,minimum height=.9cm},
terminal/.style={draw,double,rounded corners,fill=green!8,align=center,minimum width=2.6cm,minimum height=.9cm},
arr/.style={-{Latex[length=1.7mm]},thick}]
\node[state](notstarted){Not started};
\node[state,right=1.2cm of notstarted](inprogress){In progress\\percentage $<100$};
\node[terminal,right=1.2cm of inprogress](complete){Course completed\\completion reward once};
\node[state,below=1.2cm of inprogress](retry){Duplicate / retry\\no protected reward};
\draw[arr](notstarted)--(inprogress);
\draw[arr](inprogress)edge[loop above]node{new activity}(inprogress);
\draw[arr](inprogress)--(complete);
\draw[arr](complete)edge[loop above]node{repeat}(complete);
\draw[arr](inprogress)--(retry);
\draw[arr](retry)--(inprogress);
\end{tikzpicture}}
\caption{Learning-progress state and retry behavior}\label{fig:progress-state}
\end{figure}}

\newcommand{\cacheflowdiagram}{
\begin{figure}[H]\centering
\resizebox{0.94\textwidth}{!}{
\begin{tikzpicture}[
box/.style={draw,rounded corners,align=center,minimum width=2.8cm,minimum height=.9cm},
decision/.style={draw,diamond,aspect=2,fill=orange!9,align=center},
arr/.style={-{Latex[length=1.7mm]},thick}]
\node[box,fill=teal!7](request){API read request};
\node[box,fill=blue!7,right=1cm of request](redis){Redis cache\\domain TTL};
\node[decision,right=1cm of redis](hit){Hit?};
\node[box,fill=green!8,right=1cm of hit,yshift=.9cm](return){Return cached DTO};
\node[box,fill=blue!7,right=1cm of hit,yshift=-.9cm](db){Query PostgreSQL};
\node[box,fill=teal!7,right=1cm of db](populate){Map result + cache};
\node[box,fill=red!6,below=1.25cm of redis](mutation){Content mutation\\evict affected caches};
\draw[arr](request)--(redis);\draw[arr](redis)--(hit);
\draw[arr](hit)--node[above,sloped]{yes}(return);
\draw[arr](hit)--node[below,sloped]{no}(db);\draw[arr](db)--(populate);
\draw[arr,bend left=22](populate.north)to(return.south);
\draw[arr](mutation)--(redis);
\end{tikzpicture}}
\caption{Cache-aside read and invalidation flow}\label{fig:cache-flow}
\end{figure}}

\newcommand{\embeddingsequencediagram}{
\begin{figure}[H]\centering
\resizebox{\textwidth}{!}{
\begin{tikzpicture}[x=2.55cm,y=1.02cm,arr/.style={-{Latex[length=1.6mm]},thick}]
\foreach \x/\name in {0/Admin,1/API,2/PostgreSQL,3/RabbitMQ,4/Consumer}
  {\node[draw,rounded corners,fill=teal!7,minimum width=1.9cm] at (\x,0) {\name};
   \draw[dashed] (\x,-.35)--(\x,-7.3);}
\draw[arr](0,-1)--node[above]{1. update content}(1,-1);
\draw[arr](1,-1.8)--node[above]{2. transaction write}(2,-1.8);
\draw[arr](2,-2.5)--node[above]{3. commit success}(1,-2.5);
\draw[arr](1,-3.2)--node[above]{4. publish after commit}(3,-3.2);
\draw[arr](1,-4)--node[above]{5. HTTP success}(0,-4);
\draw[arr](3,-4.8)--node[above]{6. deliver event}(4,-4.8);
\draw[arr](4,-5.6)--node[above]{7. read content}(2,-5.6);
\draw[arr](4,-6.3)--node[above]{8. write MiniLM vector}(2,-6.3);
\draw[arr,dashed](4,-7)--node[above]{9. invalidate semantic caches}(1,-7);
\end{tikzpicture}}
\caption{Post-commit content embedding sequence}\label{fig:embedding-sequence}
\end{figure}}

\newcommand{\monitoringflowdiagram}{
\begin{figure}[H]\centering
\resizebox{0.96\textwidth}{!}{
\begin{tikzpicture}[
source/.style={draw,rounded corners,fill=blue!6,align=center,minimum width=2.6cm,minimum height=.85cm},
monitor/.style={draw,rounded corners,fill=orange!10,align=center,minimum width=2.8cm,minimum height=.9cm},
arr/.style={-{Latex[length=1.7mm]},thick}]
\node[source](api){API / Actuator\\custom + JVM metrics};
\node[source,below=.55cm of api](deps){PostgreSQL, Redis\\RabbitMQ exporters};
\node[source,below=.55cm of deps](probe){Blackbox probe\\main API listener};
\node[monitor,right=1.7cm of deps](prom){Prometheus\\scrape, store, rules};
\node[monitor,right=1.5cm of prom,yshift=.9cm](grafana){Grafana\\3 dashboards};
\node[monitor,right=1.5cm of prom,yshift=-.9cm](alert){Alertmanager\\group and route};
\node[monitor,right=1.3cm of alert](operator){Operator\\local inspection};
\draw[arr](api)--(prom);\draw[arr](deps)--(prom);\draw[arr](probe)--(prom);
\draw[arr](prom)--(grafana);\draw[arr](prom)--(alert);\draw[arr,dashed](alert)--node[above]{UI only}(operator);
\draw[arr,dashed,bend left=28](operator.north)to(grafana.east);
\end{tikzpicture}}
\caption{Metrics, visualization, and alerting flow}\label{fig:monitoring-flow}
\end{figure}}
